\section{Introduction}
\label{sec:intro}

% Rubric: problem understanding, project idea and design rationale (8)
The task is simulated in Webots~\cite{michel2004webots} with an e-puck robot~\cite{mondada2009epuck}. \TODO{The problem: an e-puck in a known arena with eight observation stations; one station displays the mission target. Success = stopping within 0.20\,m of the correct observation position, no collision, within 4:00 simulation time.}

\TODO{What makes it hard: target identity is only available visually (no world-file or texture access); low camera sees only the lower half of each poster; several stations are in shadow; robot arrives with heading error; narrow corridors.}

\TODO{Project idea in one paragraph: map-based navigation to candidate stations + visual verification at each station + two-layer safety monitor.}

\paragraph{Contributions.}
\begin{itemize}[nosep]
  \item \TODO{Perception: DINOv3 ViT-S/16 fine-tuned on eight reference images with view-matched augmentation; 32\,ms/frame on CPU.}
  \item \TODO{Planning: grid-based shortest-path planning and nearest-station visiting order.}
  \item \textbf{Safety:} A two-layer safety monitor combining per-step directional proximity monitoring with progress-based stuck detection for contacts missed by the proximity sensors, with replanning after each safety event. 
  The front thresholds were empirically calibrated, improving stopping clearance in a 45$^\circ$ approach from 3.6\,cm to 5.2\,cm, and the stuck detector was validated through fault-injection tests. (\cref{sec:safety})
  \item \TODO{Evaluation: three student missions plus a 96-run Monte Carlo benchmark with an independent referee.}
\end{itemize}
